\documentclass[10pt,letterpaper]{article}
\usepackage[margin=0.78in,headheight=17pt]{geometry}
\usepackage[T1]{fontenc}
\usepackage[utf8]{inputenc}
\usepackage{lmodern,microtype,parskip}
\usepackage{amsmath,amssymb}
\usepackage{xcolor,booktabs,longtable,tabularx,array,enumitem}
\usepackage[most]{tcolorbox}
\usepackage{fancyhdr,lastpage,titlesec,hyperref,listings}
\usepackage{tikz}
\usetikzlibrary{arrows.meta,positioning,shapes.geometric}
\definecolor{navy}{HTML}{17324D}
\definecolor{bluegray}{HTML}{43657D}
\definecolor{lightblue}{HTML}{EAF2F8}
\definecolor{linegray}{HTML}{CBD5DF}
\definecolor{softgray}{HTML}{F3F5F7}
\definecolor{accent}{HTML}{8E3B46}
\hypersetup{colorlinks=true,linkcolor=navy,urlcolor=bluegray,pdfborder={0 0 0}}
\urlstyle{same}
\titleformat{\section}{\Large\bfseries\color{navy}}{\thesection}{0.6em}{}
\titleformat{\subsection}{\large\bfseries\color{bluegray}}{\thesubsection}{0.5em}{}
\setlist[itemize]{leftmargin=*,itemsep=3pt,topsep=3pt}
\setlist[enumerate]{leftmargin=*,itemsep=3pt,topsep=3pt}
\setlength{\emergencystretch}{3em}
\renewcommand{\arraystretch}{1.22}
\setlength{\tabcolsep}{4pt}
\pagestyle{fancy}\fancyhf{}
\lhead{Secure Programming \& SAST Engineering}
\rhead{Engineering Documentation Suite}
\cfoot{\thepage\ / \pageref{LastPage}}
\newcolumntype{Y}{>{\raggedright\arraybackslash}X}
\newcommand{\checkbox}{\(\square\)}
\newcommand{\blankline}{\noindent\rule{\linewidth}{0.35pt}\par\vspace{4pt}}
\newcommand{\writearea}[1]{\par\begin{tcolorbox}[colback=white,colframe=linegray,boxrule=0.5pt,height=#1,sharp corners]\mbox{}\end{tcolorbox}\par}
\newtcolorbox{focusbox}[1]{colback=lightblue,colframe=bluegray,title=\textbf{#1},fonttitle=\bfseries,breakable}
\newtcolorbox{warnbox}[1]{colback=softgray,colframe=accent,title=\textbf{#1},fonttitle=\bfseries,breakable}
\lstset{basicstyle=\ttfamily\footnotesize,breaklines=true,frame=single,rulecolor=\color{linegray},backgroundcolor=\color{softgray},showstringspaces=false,columns=fullflexible,keepspaces=true,upquote=true}

\begin{document}
\begin{titlepage}
\vspace*{1.0in}
{\color{bluegray}\Large Hands-on laboratory}\\[1.0em]
{\color{navy}\Huge\bfseries Python Secure Coding Laboratory}\\[0.7em]
{\Large Paths, subprocess interfaces, SQL, privacy-safe logging and SAST}\\[2.0em]
\rule{\linewidth}{1.0pt}
\vfill
\textbf{Source discipline:} Book-derived foundations are identified separately from current engineering extensions.\\
\textbf{Edition:} 29 October 2026
\end{titlepage}
\tableofcontents
\clearpage
\section{Objective and environment}
Implement secure filesystem and subprocess boundaries in Python 3.11+, using a local disposable directory. Compare secure parameterized database calls to string interpolation and trace how request data reaches operating-system APIs.
\section{Exercise P1: Path containment}
Use \texttt{examples/python/safe\_paths.py}. Model the risk: a caller can influence a relative filename while a trusted root directory is fixed. Reject absolute paths, parent traversal that escapes the root and attempts to use an unexpected file type. Use \texttt{Path.resolve()} and \texttt{Path.is\_relative\_to()} for lexical/canonical containment, but document symlink races if hostile actors can mutate the filesystem.
\section{Exercise P2: Argument-vector execution}
\textbf{Anti-pattern for inspection only:}
\begin{lstlisting}[language=Python]
subprocess.run("grep " + pattern + " " + name, shell=True)
\end{lstlisting}
\textbf{Safer starting point for a trusted utility and validated file:}
\begin{lstlisting}[language=Python]
subprocess.run(["grep", "--", pattern, str(path)],
               check=True, shell=False)
\end{lstlisting}
\textbf{Caveat:} even without shell execution, command-specific option processing and untrusted paths require deliberate handling; prefer in-process APIs where practical.
\section{Exercise P3: SQL and logging}
Use DB-API parameter binding (placeholder style is driver-specific). Reject any implementation that builds SQL clauses from user text directly. Provide structured exception logging that does not serialize tokens, headers or query values. Distinguish a credential secret scanner from a general-purpose SAST dataflow rule.
\section{Exercise P4: Reproducible scanning and tests}
\begin{lstlisting}[language=bash]
python3 -m unittest discover -s examples/python -p 'test_*.py'
# In a controlled environment with a supported CLI:
semgrep scan --config auto examples/python
\end{lstlisting}
Record analyzer version, rule pack, target tree, excluded paths and findings. Do not treat a zero-warning scan as complete proof; add tests for absolute, relative, traversal and Unicode path cases.
\section{Required submission}
Before/after snippets, tests, sanitized scan output, residual race assumptions and a completed triage record. Pass: safe path rejection tests, no attacker-built shell command, SQL parameter binding, and reproducibility metadata.
\end{document}
